(1) 定义一维快速度 \(\chi\) 使 \(v/c=\tanh\chi\)，对应四速度 \(u^\mu=c(\cosh\chi,\sinh\chi)\)。四加速度的闵氏模为 \(a=c\,d\chi/d\tau\)，固有加速度恒定且初速为零，所以 \(\chi(\tau)=a\tau/c\)。这种参数化自动满足 \(|v|<c\)。
(2) 由 \(dt/d\tau=\cosh(a\tau/c)\) 与 \(dx/d\tau=c\sinh(a\tau/c)\) 积分得 \(t=(c/a)\sinh(a\tau/c)\)、\(x=(c^2/a)[\cosh(a\tau/c)-1]\)。消去 \(\tau\) 得 \((x+c^2/a)^2-c^2t^2=(c^2/a)^2\)，且 \(v(t)=at/\sqrt{1+(at/c)^2}\)，速度渐近 \(c\) 但不越过。
(3) 粒子总能量 \(E=mc^2\cosh\chi\)，动能 \(K=mc^2(\cosh\chi-1)\)，动量 \(p=mc\sinh\chi\)。固有时增长使 \(K,p\) 无上界，但速度 \(v=c\tanh\chi\) 仅渐近 \(c\)，因此不能用牛顿形式 \(v=a\tau\) 处理长期加速。
(4) 在瞬间共动参考系中向后发光的四波矢为 \((\omega'/c,-\omega'/c)\)。用向右快速度 \(\chi\) 的洛伦兹变换到实验系，光频率为 \(\omega=\gamma\omega'(1-\beta)=\omega'e^{-\chi}\)，故 \(\nu_{\rm lab}/\nu_{\rm ship}=\exp(-a\tau/c)\)。该式是两惯性系瞬时测频的多普勒因子，不包含之后传播路径或接收器引力红移。
